\begin{pastexam}[label={px:ccu-110-swa-11}]{中正資工所甲組 110 學年度 第 14 題}
（程式）

(10\%) Please write a function to find a key in a BST (Binary Search Tree). Return the address of
the matched node if there is a match. Otherwise, return NULL;\\
The data structure for a tree node contains a key and two pointers. The key is a string. You need
to declare the data structure of the tree node before the function code.
\end{pastexam}

\begin{pastsoln}
做法就是資結書虛擬碼~9.2 的迴圈版搜尋（遞迴版是資結書虛擬碼~9.1）：從根開始，要找的鍵比目前
節點的鍵小就往左、大就往右，相等就找到，走到空指標就表示不在樹裡。 鍵是字串，「比大小」要用
\texttt{strcmp}，依字典順序比。

\begin{code}{c}{字串鍵的 BST 搜尋}
#include <string.h>

typedef struct node {
    char *key;            /* 以 '\0' 結尾的字串 */
    struct node *left;    /* 鍵比較小的子樹 */
    struct node *right;   /* 鍵比較大的子樹 */
} Node;

Node *bst_search(Node *root, const char *key)
{
    Node *x = root;
    while (x != NULL) {
        int c = strcmp(key, x->key);
        if (c == 0)
            return x;              /* 找到了 */
        x = (c < 0) ? x->left : x->right;
    }
    return NULL;                   /* 走到空的子樹：不在樹裡 */
}
\end{code}

\textbf{幾個容易錯的地方}：
\begin{itemize}
\item 不能寫 \texttt{key == x->key}。 兩邊都是 \texttt{char *}，\texttt{==} 比的是兩個指標的位址，
內容相同、存在不同地方的字串會被當成不相等。 \texttt{<}、\texttt{>} 也一樣，要用 \texttt{strcmp}
的回傳值：負數表示 \texttt{key} 排在前面（往左），正數往右，$0$ 是相等。
\item \texttt{strcmp} 每個迴圈只呼叫一次，結果存起來用兩次，不必比兩遍。
\item 樹是空的（\texttt{root} 是 \texttt{NULL}）時迴圈一次都不做，直接回傳 \texttt{NULL}。
\item 結構裡的 \texttt{key} 也可以宣告成固定長度的陣列，例如 \texttt{char key[32];}，搜尋函式不用改。
\end{itemize}

\textbf{正確性}。 迴圈不變量：「如果 \texttt{key} 在樹裡，它在以 \texttt{x} 為根的子樹裡。」 一開始
\texttt{x} 是根，成立。 \texttt{key} 比 \texttt{x->key} 小時，由 BST 性質，右子樹的鍵都比
\texttt{x->key} 大，不可能是 \texttt{key}，所以 \texttt{key} 若在就在左子樹，往左走後仍成立；比較大時
同理。 回傳 \texttt{x} 時 \texttt{strcmp} 是 $0$，確實找到；\texttt{x} 變成 \texttt{NULL} 時子樹是空的，
由不變量 \texttt{key} 不在樹裡，回傳 \texttt{NULL} 正確。

\textbf{時間}。 每一圈往下走一層，最多走 $h + 1$ 圈，$h$ 是樹高；每圈一次 \texttt{strcmp}，最多比到
較短的字串結束，設字串長度不超過 $L$，總時間 $O(hL)$。 鍵長是常數時就是資結書第~9.1~節的 $O(h)$。
\end{pastsoln}
